\documentclass[aspectratio=169]{beamer}
\usetheme{Madrid}

\usepackage{amsmath}
\usepackage{amssymb}
\usepackage{booktabs}
\usepackage{graphicx}
\usepackage[strings]{underscore}

\title{\texttt{s\_lolr}: Model and Simulation Results}
\author{}
\date{\today}

\begin{document}

\begin{frame}
    \titlepage
\end{frame}

\begin{frame}{Model}
\small
\[
V^{ND}(b,l,g,s,\varepsilon)=
\max_{b',l'}
\left\{
u(c^{ND})+\beta g^{1-\gamma}\,
\mathbb{E}\!\left[\max\{V^{ND}(b',l',g',s',\varepsilon'),V^D(b',l',g',s',\varepsilon')\}\mid g,s\right]
\right\}
\]
\[
c^{ND}=y(g,\varepsilon)-b-l+n(b',l',g,s)+n_l^{ND}(l',g)
\]

\[
V^D(b,l,g,s,\varepsilon)=
\max_{l'}
\left\{
u(c^D)+\beta g^{1-\gamma}\,
\mathbb{E}\!\left[\theta V^D\!\left(\frac{b}{g},l',g',s',\varepsilon'\right)
+(1-\theta)\max\left\{V^{ND}\!\left(\kappa\frac{b}{g},l',g',s',\varepsilon'\right),V^D\!\left(\frac{b}{g},l',g',s',\varepsilon'\right)\right\}\middle| g,s\right]
\right\}
\]
\[
c^D=\phi_g\,y(g,\varepsilon)-l+n_l^{D}(l',g)
\]
\[
n_l^{ND}(l',g)=\frac{g\,l'}{R_l^{ND}},
\qquad
n_l^{D}(l',g)=\frac{g\,l'}{R_l^{D}}
\]
\end{frame}

\begin{frame}{Prices and Schedules}
\small
\[
d(b,l,g,s,\varepsilon)=\mathbf{1}\left\{V^D(b,l,g,s,\varepsilon)>V^{ND}(b,l,g,s,\varepsilon)\right\}
\]
\[
Q(b,l,g,s,\varepsilon)=(1-d)+d\,X(b,l,g,s,\varepsilon)
\]
\[
X(b,l,g,s,\varepsilon)=\frac{1}{1+r^\ast}\,
\mathbb{E}\!\left[
\theta X'
+(1-\theta)\left((1-e')X'+e'\kappa Q'\right)
\middle| g,s
\right]
\]
\[
e'(b',l',g',s',\varepsilon')=
\mathbf{1}\left\{
V^{ND}(\kappa b',l',g',s',\varepsilon')\geq V^D(b',l',g',s',\varepsilon')
\right\}
\]
\[
R(b',l',g,s)=\frac{1+r^\ast}{\mathbb{E}[Q' \mid g,s]},
\qquad
n(b',l',g,s)=\frac{g\,b'}{R(b',l',g,s)}
\]
\[
p^{def}(b',l',g,s)=\mathbb{E}[d' \mid g,s]
\]
\[
p^{def}(b',l',g,s)\leq \bar p_b,
\qquad
n(b',l',g,s)\leq \bar n_g,
\qquad
l' \leq \bar l^{ND}, \bar l^{D}
\]
\end{frame}

\begin{frame}[shrink=18]{Simulation Table}
\hypertarget{main-sim-table}{}
\scriptsize
\centering
\begingroup
\setlength{\tabcolsep}{4pt}
\renewcommand{\arraystretch}{0.92}
\input{../../1period_s_lolr_v2/result/s_lolr_v2_moments_table.tex}
\endgroup

\medskip
\hyperlink{transition-table-appendix}{\beamergotobutton{Transition Table}}
\hspace{0.6em}
\hyperlink{all-moments-table-appendix}{\beamergotobutton{test moments}}
\hspace{0.6em}
\hyperlink{all-transition-table-appendix}{\beamergotobutton{test Transition}}
\hspace{0.6em}
\hyperlink{sim-var-appendix}{\beamergotobutton{Simulation Variables}}
\end{frame}

\begin{frame}{Selected Private Schedule: Baseline vs V2}
\begin{columns}[T,onlytextwidth]
\column{0.33\textwidth}
\centering
\includegraphics[width=\linewidth]{../../1period_s_lolr/result/s_lolr_low_growth_Rn_eqm_lprime.png}

{\scriptsize Fixed LOLR borrowing}

\column{0.33\textwidth}
\centering
\includegraphics[width=\linewidth]{../../1period_s_lolr_v2/result/s_lolr_v2_low_growth_selected_schedule.png}

{\scriptsize State-dependent LOLR borrowing}

\column{0.34\textwidth}
\scriptsize
\begin{itemize}
    \item Conditional private schedule at $l'=0$
    \item Low growth; vary $b'$ and plot $\left(n(b',0,g,s),R(b',0,g,s)\right)$
    \item Solid: bad $s$ \quad Dashed: good $s$
    \item Left colors compare fixed-$\Delta$ models
    \item Right compares bad/good $s$ within one state-dependent-$\Delta$ model
\end{itemize}
\end{columns}
\end{frame}

\begin{frame}{Private Borrowing Policy: Baseline vs V2}
\begin{columns}[T,onlytextwidth]
\column{0.33\textwidth}
\centering
\includegraphics[width=\linewidth]{../../1period_s_lolr/result/s_lolr_low_growth_nb_eqm_lprime.png}

{\scriptsize Fixed-$\Delta$ models}

\column{0.33\textwidth}
\centering
\includegraphics[width=\linewidth]{../../1period_s_lolr_v2/result/s_lolr_v2_low_growth_private_policy.png}

{\scriptsize State-dependent $\Delta$}

\column{0.34\textwidth}
\scriptsize
\textbf{}
\begin{itemize}
    \item Plot
    \[
    n(b)=n\!\left(b^\ast(b),l^\ast(b),g,s\right)
    \]
    with current $l=0$ and low growth.
\end{itemize}
\textbf{Left: baseline}
\begin{itemize}
    \item Separate fixed-$\Delta$ models; colors index $\Delta$.
    \item Solid/dashed split bad and good $s$.
\end{itemize}
\textbf{Right: V2}
\begin{itemize}
    \item One model with state-dependent LOLR cap.
    \item Difference between solid and dashed now comes directly from changing the available LOLR cap across sunspots.
\end{itemize}
\end{columns}
\end{frame}

\begin{frame}{LOLR Borrowing Policy: Baseline vs V2}
\begin{columns}[T,onlytextwidth]
\column{0.33\textwidth}
\centering
\includegraphics[width=\linewidth]{../../1period_s_lolr/result/s_lolr_low_growth_nl_policy.png}

{\scriptsize Fixed-$\Delta$ models}

\column{0.33\textwidth}
\centering
\includegraphics[width=\linewidth]{../../1period_s_lolr_v2/result/s_lolr_v2_low_growth_lolr_policy.png}

{\scriptsize State-dependent $\Delta$}

\column{0.34\textwidth}
\scriptsize
\textbf{}
\begin{itemize}
    \item Left subpanel:
    \[
    n_l^{ND}(b)=\frac{g\,l^\ast_{ND}(b)}{R_l^{ND}}
    \]
    \item Right subpanel:
    \[
    n_l^{D}(b)=\frac{g\,l^\ast_{D}(b)}{R_l^{D}}
    \]
    \item Both use the full current $b$ grid with current $l=0$.
\end{itemize}
\textbf{Difference}
\begin{itemize}
    \item Baseline colors compare fixed-$\Delta$ models.
    \item V2 compares bad and good $s$ within one state-dependent-$\Delta$ model, so the sunspot itself changes LOLR access in repayment.
\end{itemize}
\end{columns}
\end{frame}

\begin{frame}{Default Regions: Baseline vs V2}
\begin{columns}[T,onlytextwidth]
\column{0.33\textwidth}
\centering
\includegraphics[width=\linewidth]{../../1period_s_lolr/result/s_lolr_low_growth_default_policy.png}

{\scriptsize Fixed-$\Delta$ models}

\column{0.33\textwidth}
\centering
\includegraphics[width=\linewidth]{../../1period_s_lolr_v2/result/s_lolr_v2_low_growth_default_policy.png}

{\scriptsize State-dependent $\Delta$}

\column{0.34\textwidth}
\scriptsize
\textbf{}
\begin{itemize}
    \item Plot
    \[
    d(b,0,g,s,\varepsilon)=\mathbf{1}\{V^D>V^{ND}\}
    \]
    over the full current $b$ grid at low growth.
\end{itemize}
\textbf{Left: baseline}
\begin{itemize}
    \item Different colors are different constant-$\Delta$ models.
\end{itemize}
\textbf{Right: V2}
\begin{itemize}
    \item Same model in both lines.
    \item Bad $s$ has higher repayment LOLR cap than good $s$, so any shift is the direct within-model effect of state-dependent access.
\end{itemize}
\end{columns}
\end{frame}

\begin{frame}{State-Dependent LOLR: Distribution of Private Issuance}
\begin{columns}[T,onlytextwidth]
\column{0.7\textwidth}
\centering
\includegraphics[width=\linewidth]{../../1period_s_lolr_v2/result/s_lolr_v2_issuance_distribution.png}

\column{0.3\textwidth}
\scriptsize
\begin{itemize}
    \item Histogram of simulated $n_t$
    \item Low-growth panel on the left
    \item High-growth panel on the right
    \item No-default observations only
    \item Bars show number of observations, not density
\end{itemize}
\end{columns}
\end{frame}

\begin{frame}{State-Dependent LOLR: $g_H \rightarrow g_L$ Distributions}
\begin{columns}[T,onlytextwidth]
\column{0.7\textwidth}
\centering
\includegraphics[width=\linewidth]{../../1period_s_lolr_v2/result/s_lolr_v2_transition_distributions.png}

\column{0.3\textwidth}
\scriptsize
\begin{itemize}
    \item Top row: carried debt service $b_t$
    \item Bottom row: current LOLR issuance $n_{l,t}$
    \item Left column: all $g_H \rightarrow g_L$ transitions
    \item Right column: $g_H \rightarrow g_L$ transitions with no default
    \item Bars show number of observations, not density
\end{itemize}
\end{columns}
\end{frame}

\appendix

\begin{frame}{Appendix: Simulation Variables}
\hypertarget{sim-var-appendix}{}
\small
\[
f_t =
\underbrace{g_t}_{\substack{\text{scale} \\ \text{not discounting}}}
\quad
\underbrace{b_t'}_{\substack{\text{principal payment} \\ \text{based on budget constraint}}}
\]
\[
q_t = \frac{1}{R_t} = \frac{E_t[Q_{t+1}]}{1+r^*}
\]
\[
q_t b_t = \frac{b_t}{R_t}
\]
\[
n_t = q_t f_t = q_t g_t b_t' = \frac{g_t\, b_t'}{R_t}
\]

\medskip
\begin{itemize}
    \item $f_t$ is the face value of newly issued private debt.
    \item $n_t$ is private proceeds today from newly issued debt.
    \item $q_t b_t$ is the market value of inherited private debt.
\end{itemize}

\medskip
\hyperlink{main-sim-table}{\beamerreturnbutton{Back to Main Table}}
\end{frame}

\begin{frame}[shrink=10]{Appendix: Default Transition Table}
\hypertarget{transition-table-appendix}{}
\tiny
\centering
\begingroup
\setlength{\tabcolsep}{2pt}
\renewcommand{\arraystretch}{0.85}
\input{../../1period_s_lolr_v2/result/s_lolr_v2_transition_table.tex}
\endgroup

\medskip
\hyperlink{main-sim-table}{\beamerreturnbutton{Back to Main Table}}
\end{frame}

\begin{frame}[shrink=18]{Appendix: All Moments Table}
\hypertarget{all-moments-table-appendix}{}
\tiny
\centering
\begingroup
\setlength{\tabcolsep}{2pt}
\renewcommand{\arraystretch}{0.82}
\input{../../1period_s_lolr_v2/result/s_lolr_v2_all_moments_table.tex}
\endgroup

\medskip
\hyperlink{main-sim-table}{\beamerreturnbutton{Back to Main Table}}
\end{frame}

\begin{frame}[shrink=12]{Appendix: All Transition Table}
\hypertarget{all-transition-table-appendix}{}
\tiny
\centering
\begingroup
\setlength{\tabcolsep}{2pt}
\renewcommand{\arraystretch}{0.82}
\input{../../1period_s_lolr_v2/result/s_lolr_v2_all_transition_table.tex}
\endgroup

\medskip
\hyperlink{main-sim-table}{\beamerreturnbutton{Back to Main Table}}
\end{frame}

\end{document}
